% ─────────────────────────────────────────────────────────────────────────────
% CAPÍTULO 2 — Pipeline de Voz e Inventario de Servicios
% ─────────────────────────────────────────────────────────────────────────────
\chapter{Pipeline de Voz e Inventario de Servicios}

\section{Pipeline End-to-End de Voz}

El objetivo central del Hermes Stack es procesar voz en tiempo real con
una latencia total menor a 500 ms. El pipeline consta de cinco etapas
secuenciales:

\begin{figure}[h]
\centering
\begin{tikzpicture}[
  node distance=0cm and 0.3cm,
  stage/.style={draw=hermesblue,fill=hermeslight,rounded corners=5pt,
                minimum width=2.5cm,minimum height=1.2cm,
                font=\sffamily\small\bfseries,align=center},
  lat/.style={font=\sffamily\tiny\color{greenok},below=2pt},
  arr/.style={->,>=Stealth,very thick,hermesblue}
]
\node[stage] (mic)    {Micrófono\\\small Audio PCM};
\node[stage,right=0.5cm of mic]     (wsp)  {Whisper\\STT};
\node[stage,right=0.5cm of wsp]     (hrm)  {Hermes\\Agent};
\node[stage,right=0.5cm of hrm]     (llm)  {LiteLLM\\→ LLM};
\node[stage,right=0.5cm of llm]     (tts)  {ElevenLabs\\Flash v2.5};
\node[stage,right=0.5cm of tts,fill=hermesblue!30] (out) {Audio\\Salida};

\draw[arr] (mic) -- (wsp);
\draw[arr] (wsp) -- (hrm);
\draw[arr] (hrm) -- (llm);
\draw[arr] (llm) -- (tts);
\draw[arr] (tts) -- (out);

\node[lat] at ($(mic)!0.5!(wsp)+(0,-0.7)$) {$\sim$80--120 ms};
\node[lat] at ($(wsp)!0.5!(hrm)+(0,-0.7)$) {$\sim$5 ms};
\node[lat] at ($(hrm)!0.5!(llm)+(0,-0.7)$) {$\sim$100--180 ms};
\node[lat] at ($(llm)!0.5!(tts)+(0,-0.7)$) {$\sim$150--220 ms};
\end{tikzpicture}
\caption{Pipeline de voz end-to-end (latencia total objetivo: $<$500 ms)}
\end{figure}

\subsection{Latencias por etapa}

\begin{table}[h]
\centering
\caption{Desglose de latencias en el pipeline de voz}
\begin{tabularx}{\textwidth}{llXl}
\toprule
\textbf{Etapa} & \textbf{Servicio} & \textbf{Descripción} & \textbf{Latencia} \\
\midrule
1. STT        & Whisper (base) & Transcripción audio $\to$ texto      & 80--120 ms \\
2. Routing    & Hermes Agent   & Parsing + selección de skill          & 3--8 ms    \\
3. LLM        & LiteLLM        & Primer token (TTFB) desde OpenRouter  & 100--180 ms \\
4. TTS        & ElevenLabs WS  & Streaming Flash v2.5 WebSocket        & 150--220 ms \\
5. Red/Buffer & VPS $\to$ cliente & Transmisión y buffering de audio   & 10--30 ms  \\
\midrule
\textbf{Total E2E} & & & \textbf{343--558 ms} \\
\bottomrule
\end{tabularx}
\end{table}

\subsection{Optimizaciones de latencia}

\begin{itemize}
  \item \textbf{WebSocket persistente con ElevenLabs:} Se reutiliza la
    conexión WS entre peticiones para eliminar el handshake TLS ($\sim$50 ms
    de ahorro por turno).
  \item \textbf{Streaming de tokens LLM:} LiteLLM reenvía tokens en
    cuanto llegan; Hermes inicia la llamada TTS con el primer fragmento
    coherente, no al finalizar la respuesta completa.
  \item \textbf{Modelo Whisper base:} Elegido sobre \code{small} o
    \code{medium} para garantizar transcripción en $<$120 ms en CPU.
  \item \textbf{Redis como caché semántico:} LiteLLM usa Redis para
    respuestas exactamente repetidas (saludos, comandos frecuentes).
\end{itemize}

\subsection{Política de fallback en LiteLLM}

\begin{lstlisting}[language=yaml,caption=Fragmento de config/litellm.yaml — fallback]
router_settings:
  routing_strategy: latency-based-routing
  fallbacks:
    - model: openrouter/meta-llama/llama-3.1-70b-instruct
      fallback_models:
        - gemini/gemini-flash-1.5
        - openrouter/mistralai/mistral-7b-instruct
  timeout: 10
  num_retries: 2
\end{lstlisting}

\section{Inventario Detallado de Servicios}

\subsection{LiteLLM Router (:4000)}

\service{LiteLLM} actúa como proxy inteligente entre Hermes y los
proveedores de LLM externos. Centraliza la autenticación, el enrutamiento
basado en latencia y el caché Redis.

\begin{lstlisting}[language=yaml,caption=docker-compose.yml — servicio litellm]
litellm:
  image: ghcr.io/berriai/litellm:main-latest
  environment:
    LITELLM_MASTER_KEY: ${LITELLM_MASTER_KEY}
    REDIS_HOST: redis
    OPENROUTER_API_KEY: ${OPENROUTER_API_KEY}
    GEMINI_API_KEY: ${GEMINI_API_KEY}
  volumes:
    - ./config/litellm.yaml:/app/config.yaml:ro
  ports:
    - "127.0.0.1:4000:4000"
\end{lstlisting}

\begin{table}[h]
\centering
\caption{Variables de entorno requeridas por LiteLLM}
\begin{tabularx}{\textwidth}{lX}
\toprule
\textbf{Variable} & \textbf{Descripción} \\
\midrule
\envvar{LITELLM\_MASTER\_KEY} & Clave maestra para autenticar peticiones al proxy \\
\envvar{OPENROUTER\_API\_KEY} & Clave de OpenRouter para acceso a múltiples LLMs \\
\envvar{GEMINI\_API\_KEY}     & Clave de Google Gemini como fallback primario \\
\bottomrule
\end{tabularx}
\end{table}

\subsection{Hermes Agent (:8080)}

\service{Hermes Agent} es el núcleo del sistema. Implementado en Python,
expone una API HTTP/WebSocket que orquesta el pipeline completo de voz.

Responsabilidades principales:
\begin{itemize}
  \item Recibir audio del cliente vía WebSocket
  \item Llamar a Whisper STT para transcripción
  \item Enrutar la petición a LiteLLM con el prompt de sistema adecuado
  \item Transmitir la respuesta de texto a ElevenLabs Flash v2.5
  \item Hacer streaming del audio de vuelta al cliente
  \item Exponer métricas Prometheus en \code{/metrics}
\end{itemize}

\begin{lstlisting}[language=python,caption=Fragmento ilustrativo — pipeline de voz]
async def voice_pipeline(ws: WebSocket):
    audio_bytes = await ws.recv()
    transcript  = await whisper.transcribe(audio_bytes)
    llm_stream  = litellm.stream(transcript)
    async for chunk in elevenlabs.tts_stream(llm_stream):
        await ws.send_bytes(chunk)
\end{lstlisting}

\subsection{Whisper STT (:9000)}

\service{Whisper STT} corre el modelo \code{base} de OpenAI Whisper
localmente, garantizando privacidad total de audio.

\begin{infobox}[Configuración recomendada]
Usar \code{ASR\_MODEL: base} para latencia $<$120 ms en CPU de 2 vCPU.
Para CPU de 4+ vCPU, \code{small} puede usarse con latencia de $\sim$200 ms
pero mejor precisión en acentos.
\end{infobox}

\subsection{Redis Cache}

\service{Redis} 7 Alpine corre con capacidades mínimas:

\begin{lstlisting}[language=yaml,caption=Redis hardening]
redis:
  image: redis:7-alpine
  cap_drop: [ALL]
  cap_add: [CHOWN, SETGID, SETUID]
  tmpfs:
    - /tmp:size=32M,noexec,nosuid
\end{lstlisting}

\subsection{Hermes Website (:3001)}

\service{Hermes Website} es una SPA React/Vite con backend Express
que sirve la interfaz web del proyecto en \domain{docs.el80.space}.

\begin{table}[h]
\centering
\caption{Stack tecnológico del sitio web}
\begin{tabularx}{\textwidth}{lX}
\toprule
\textbf{Capa} & \textbf{Tecnología} \\
\midrule
Frontend & React 18 + Vite + TypeScript + Tailwind CSS \\
Backend  & Express (Node 22) — proxy y endpoint \code{/api/tree} \\
Build    & Docker multi-stage (builder + runner), imagen Alpine \\
Assets   & Framer Motion para animaciones, Heroicons para iconos \\
\bottomrule
\end{tabularx}
\end{table}

\subsection{File Browser (:8095)}

\service{File Browser} es una interfaz web para gestión de archivos del
servidor, accesible en \domain{files.el80.space}. Monta \code{/root}
como raíz del filesystem navegable.

\begin{warnbox}[Seguridad — File Browser]
File Browser corre como root (\code{user: "0:0"}) con acceso de escritura
total a \code{/root}. Debe protegerse con contraseña fuerte en la UI y
acceso SSH-only al VPS.
\end{warnbox}
